\documentclass{article}
\usepackage{graphicx} % Required for inserting images

\title{5CL Lab 5}
\author{Isaac Chan}
\date{October 2026}

\begin{document}

\maketitle

\section{Helium}
\subsection{1.1}
Spectrum reader: 9.13
lens: 45.71

We observe bright colored lights along the spectrum, which eventually repeat once much fainter at extreme angles. This repeat (the 2nd order) basically only has yellow line apparent

We align the reader by making sure values on both sides read about the same - small sanity check

Isaac:
On the right, yellow is at 11.11cm, uncertainty .15cm. Left 11.18, uncertainty .15

Forrest:
On the right, yellow is at $11.3\pm0.1$, left: $11.3\pm0.1$

Width of the lines is on the on the order of .1cm

\subsection{1.2}
Only record on the right side

violet, green, red,

red: $12.87 \pm 0.1$ cm
green: $9.51 \pm 0.1$
violet: $8.46 $

For the grating mirror, we roughly align mirror with the ionized gas; however our entire setup is off of the optical bench and thus it becomes very hard to contrain any value

\section{Hydrogen}

We use and identical setup to before:
Red Cyan Blue/Violet

Reader: 14.35
Lens: 42.2 cm

red: $9.78\pm0.1$ 
cyan: $6.97\pm 0.1$
violet: $6.24$

\section{Mystery}

B:
red: $9.59\pm0.1$ 
orange: $9.19\pm 0.1$
yellow: $8.45$

C:
This is argon because there are a lot of lines, many of which are blue, green (in two distnct areas), and yellow-red.

\section{Light Bulbs}
Incandescent is a full, continuous specturm. Green and red seem the strongest, which lines up with literature Ive read on ocnes. The red fades very gradually, and hte violet fades faster


We put the cfl behind the slit, and then view the diffraction grating. We see a broad spectrum of light, and also see the 2nd order.

we see nie spectrum for 

We see bright red, cyan, yellow, violet lines, bluring overall for the spectrum


\end{document}
